% ============================================================
% 时代写作榜（杂志页码 68–71）
% ============================================================
\colchapter{时代写作榜}

\section*{@中国青年报：《2022开学第一课》：奋斗成就梦想}

9月1日，大型公益节目《2022开学第一课》在央视综合频道播出，以“奋斗成就梦想”为主题，将节目小课堂与社会大课堂紧密结合起来，回溯了2022年的热议话题：冬奥体育、农业生产、生态文明、航天科技等。这些领域的关键性人物在节目中现身说法：从在冬奥赛场上奋力拼搏的奥运冠军，到50年扎根土地只为培育金色种子的院士；从在荒原上种出全球面积最大的人工林场、创造“绿色奇迹”的三代务林人，到勇攀世界之巅、探源青藏高原地理气候变迁的科考队……通过“奋斗者”的经历，《2022开学第一课》传递出清晰的价值导向：青少年应该从国家、民族、时代的坐标系中找到自己的志向，把个人的小梦想融入国家民族的大梦想。

\section*{@播in：成长是一场修行}

成长是每个人终其一生的修行，在持续学习中收获超越自我的能力与心理状态，是成长的魅力所在。如果你觉得前一年的自己很蠢，那说明你这一年有了进步。保持前进的勇气、及时进行自我省察、面对未知淡然平和、不给自身设限，才会见证厚积薄发。愿你始终昂首挺胸，成为更好的自己。

\section*{@人民日报：中国科学家首次在月球发现新矿物“嫦娥石”}

2020年12月17日，嫦娥五号携带1731克月球样品返回地球，多个单位获批承担该样品的研究工作。2022年9月9日，国家航天局、国家原子能机构联合发布了针对月球样品的最新研究成果：中国科学家首次在月球上发现的新矿物，被命名为“嫦娥石”。这是我国在空间科学领域取得的一项重大成果。“嫦娥石”是一种磷酸盐矿物，呈柱状晶体，存在于月球玄武岩颗粒中，是人类在月球上发现的第六种新矿物。就此，我国成为继美国、苏联后，世界上第三个在月球发现新矿物的国家。

\clearpage
\section*{拒绝假装努力}

\subsection*{巧命题}

材料一：有的人在星巴克里放个电脑一坐就是一下午，也许只是在上网聊天；有的人在健身房“运动五分钟、自拍两小时”……这些人被称为“假装努力气氛组”。

材料二：2018年，河南一考生家长因孩子高考成绩与估分差距大，质疑答题卡被调包而进行实名举报。调查后，真相却是该考生在平时考试中大量抄袭，高考则使其原形毕露，为逃避责骂，便编造了“答题卡被调包”的谎言。

材料三：不要假装努力，因为结果不会陪你演戏。

对此，你有什么思考？请写一篇论述类文章加以阐明。

\subsection*{时材立意}

本命题材料介绍了现实中假装努力的现象及其造成的严重后果。对此，我们可以从假装努力的原因、危害和如何避免该现象的角度得出如下立意：

\textbf{①以结果为导向，拒绝“假努力”。}不论是正在成长的青少年，还是已工作的成年人，勤奋要以梦想为导向，朝着既定方向的努力才是真实而有效的。假装努力的实质，是向上的知易行难，而“行难”的背后是人的惰性，要想改变则需要有更强大的动力支撑。“少年壮志当拿云”“老骥伏枥，志在千里”，李贺和曹操告诉我们，无论什么时候都不能丧失对志向的追求。

\textbf{②勇敢自信行路，何惧他人眼光。}我们从小的学习模式就是遵循规律的上课时间、完成老师布置的作业并得到分数或者评价。有的人能从中找到自己学习的节奏并养成良好的生活习惯。而有的人一旦脱离了外在约束，就变成了只会旋转的陀螺，原地踏步。后者上学时为了老师和家长学习，进入社会后也活在他人的评价和目光中。为得到他人的正向评价而假装努力，其实是一种内心不够自信的表现。人生如果是一个舞台，那我们应当为自己而舞。

\textbf{③做一个主宰自我的人，越自律越强大。}自制力永远是成功的根本，如果自己都不能掌控自己的身体行动，变成了懒惰和不思进取的奴隶，那又如何去与世界的冷酷和生活的磨难挫折抗衡，如何享受崎岖山路过后的绝顶风景？自律是解决人生问题的主要工具，也是消除人生痛苦的重要方法。对自己要求严格一些，对自己“狠一点”，就会发现生活对你不再严苛，世界对你更加宽容，做事的结果更加美好。

（重庆市彭水中学 何烨）

\begin{quoteBox}{网友声音}
@少女兔iiilass：生活中，有些人的勤奋是表演出来的，给自己看，也给别人看：在结果不好的时候安慰自己，“不是因为我不努力，只是脑子太笨”；为了获得别人那句“你真努力”的表扬。或许，“努力”不是一场盛大的马拉松，而更像一场孤独的奋斗，不需要那么多人站在路边为你呐喊助威。不要假装努力，因为结果不会陪你演戏。
\end{quoteBox}

\clearpage
\section*{范文审阅：鲜花之外，无声苦行}
\auth{文/重庆南开中学2023级（18）班 黄一宸（得分：56分）}

生活中，有些人的勤奋是表演出来的，也许是为了给自己的失败找借口，也许是为了获得别人“你真努力”的表扬。然而，功成之路无数，皆由真正的努力铺就，行假努力之路，虽求得暂时心安，却非长久之计。毕竟，结果不会陪我们表演。

\textbf{欲功成，需踏实而全力行之。}祖逖舞剑，天未亮而勤练，日复一日，雷打不动，终成一朝猛将；车胤求学，借萤虫微光而苦读，不骄不躁，夜夜如常，终成一代学士；苏秦发奋，觉困而刺股，不眠不休，全力以赴，终成一世谋略家。古之功成者，皆以踏实努力而成事，不努力、假努力而成事者，古未有之，而况今者乎？今人于快节奏的时代中，欲功成而不欲踏实奋进者有之，更有甚者，以努力之名而行偷懒之实。欲尝成功之果而不欲受刺于努力苦行之荆棘，此举终致身败而功不成，实不可取矣！

\textbf{踏实地，万不可以欺人。}欺人易，欺事难。欺人唯得须臾之怜悯，片刻之安慰，于己身何用之有？羲之洗笔染池黑，书圣出于世；白石久观虾形，山翁出神笔。默默无闻，不张不扬，甘于平凡而苦练技艺者，真努力而成事之人也。今之学子坐于桌前，片刻难耐，看似全力奋斗，实则神游天外，动辄曰：“努力，败也，虽败犹荣。”从而换取一片“钦佩”的掌声。此真努力邪？此真浮躁也！诚然，时代之快节奏易催生“内卷”“速成”乱象，但此亦不能成“假努力”之由。成长节奏皆为自身所选，或快或慢，或急或缓，未尝有定论。何故多人为社会时钟所钳制，而不能守自我？悲夫！

守自我，勿惧路途险恶孤寂，一心向前，奋进之途定不负人！

努力之途坎坷而寂寞无闻。老聃虽为小吏，独阅周室藏书，沉淀所学，名留万古；曹雪芹虽处逆境，作《红楼梦》十余载，字字泣血传于后世。真努力而成事者，懂得坚守自我，终或易逆境，或悟至道，或成大业，或得真理。努力以成己身，付出多，得亦多，一分一毫皆能报以涌泉。

勿惧风雨，砥砺前行。努力者自有千钧之力，骄躁者难成分毫之事。切勿为时俗所挟，内心浮躁，不定内心而单以表面之努力，以求自我之安慰、他人之怜悯。沉静，踏实，笃行，坚持……长路漫漫，尽己身之力而求精进，此乃成事之正道。

\score{得分点：举例论证，使用排比的修辞手法，句式整饬，文笔富有气势。}
\score{得分点：对比论证，论述之间横贯古今，直指假努力的行为不可取。}
\score{得分点：联系生活实际，讲出当代学生假努力的现象，具有现实针对性。}
\score{得分点：“诚然”一词，以退为进，指出现代人身不由己之处，也指明其随波逐流的悲哀。}
\score{得分点：承接上文，进一步指出真正努力该如何做，有何意义。}
\score{得分点：首尾呼应，总结全文。}

\subsection*{时评链接：“假努力”是自欺欺人}

我们经常能够看到，一些办公室晚上灯火通明。这其中，有许多人是在加班加点工作，然而也有少数人把加班当“表演”，想借“灯火通明”给人留下努力工作的印象。这种加班并非在忙工作，其实是做样子的“假努力”。“假努力”危害大。人的精力是有限的，一个人如果耗费大量时间和精力用在维持“自己很努力”这个表象上，势必难以集中精力全身心投入工作，既解决不了问题，也完成不好任务，到头来只是自欺欺人，随着时间流逝，留下的只有“一地鸡毛”。

世界上许多事情，是掺不得假的。假花虽然看起来漂亮，但它不香，也结不了果。从长远来看，“假努力”不会有好结果，那些弄虚作假的自作聪明、自欺欺人，到头来只会是一场空。因为结果不会撒谎，时间也不会撒谎，“是骡子是马”，随着时间推移，最终都会见分晓。对待工作，唯有真努力、真落实，始终高标准完成，才能获得长足的进步。

（摘自《解放军报》，有删改）

\subsection*{名师态度}

使用文言文句式进行写作时，我们应当提出更高要求，使文章符合语言规范，而不是想当然地随意使用“之乎者也”，好在本文语言精练，内容充实，文思斐然。面对真假努力的话题，作者开门见山地表明自己的观点——假努力不过能求暂时心安，最终苦果却早已注定。既而，在论证过程中，文章以“欲功成”“踏实地”“守自我”三个分论点层层推进，分论点之间亦遥相呼应。此外，文章语言可圈可点，善于使用对偶句，读来朗朗上口；首尾照应，结构完整，始终紧扣中心论点展开论述，张弛有度。

——重庆南开中学 黄渝平
